

\documentclass[compacto,10pt]{pubciencias-notas}

% -- Paquetes adicionales
\usepackage{comandosEPN}

% -- Datos del libro
\autor{Departamento de Formación Básica}
\materia{Álgebra Lineal}
\nota{Semana 04 Clase 02}
\fecha{Semestre 2025-A}

\logouno[1.1cm]{Logo_EPN.eps}
\logodos[1.7cm]{Logo_DFB.eps}

\definecolor{colordef}{cmyk}{0.81,0.62,0.00,0.22}

\begin{document}

\encabezado

\vspace*{-8mm}

\setcounter{section}{1}
\section{Sistemas de ecuaciones}
\setcounter{subsection}{2}
\subsection{Matriz Inversa}

\begin{defi}\label{Sem03Clas02DefInv}
    Sea $A\in \Mat[\K]{n}{n}$. Se dice que $A$ es no singular o invertible si existe una matriz $B\in \Mat[\K]{n}{n}$ tal que 
    \[
        AB = BA = I_n
    \]
    la matriz $B$ se denomina inversa de $A$ y se la denota por $A^{-1}$. 
    
    \noindent Si no existe tal matriz, entonces se dice que $A$ es singular o no invertible. 
\end{defi}

\begin{prop}
    Si una matriz tiene inversa, la inversa es única. 
\end{prop}

\begin{comment}
\begin{proof}\hspace{0pt} Sea $A\in\Mat[\K]{n}{n}$ una matriz invertible y supongamos que $B$ y $C$ son dos inversas de $A$, por la definición \ref{Sem03Clas02DefInv} se tiene que
\[
    AB = BA = I_{n}
    \texty
    AC = CA = I_{n}.
\]
Se tiene, además, que
\begin{align*}
    B 
    &= BI_{n}\\
    &= B(AC)\\
    &= (BA)C\\
    &= I_{n}C\\
    &= C
\end{align*}
lo que muestra que $B = C$, y, por tanto, si la inversa de $A$ existe, es única.\qedhere
\end{proof}
\end{comment}

\begin{ejem}\hspace{0pt} 
Sean $A\in\Mat[\K]{n}{n}$ una matriz no singular y sean $x, b \in \Mat[\K]{n}{1}$. Si $Ax = b$, entonces $x = A^{- 1}b$.

\begin{proof}[Solución.]\hspace{0pt} Por hipótesis sobre $A$ se tiene que $A^{-1}$ existe, entonces
\begin{align*}
    Ax = b
    &\Longleftrightarrow A^{-1}(Ax) = A^{-1}b
    && \text{Multiplicando $A^{- 1}$ en ambos lados}\\
    &\Longleftrightarrow (A^{-1}A)x = A^{-1}b
    && \text{Propiedad asociativa}\\
    &\Longleftrightarrow I_{n}x = A^{-1}b
    && \text{Definición de matriz inversa}\\
    &\Longleftrightarrow x = A^{-1}b
    && \text{Multiplicación por la matriz identidad}\qedhere
\end{align*}
\end{proof}
\end{ejem}

\begin{ejer}\hspace{0pt} 
Sean $A\in\Mat[\K]{n}{n}$ una matriz no singular y sean $x, b \in \Mat[\K]{n}{1}$. Si $x^{\intercal}A^{2} + b^{\intercal}A = 0$, entonces $x^{\intercal} = - b^{\intercal}A^{- 1}$.
\end{ejer}

\begin{teo}\label{Sem03Clas02PropInv}
    Sea $A, B \in \Mat[\K]{n}{n}$. 
    \begin{enumerate}
    \item[$a)$] 
        Si $A$ es una matriz no singular, entonces $A^{-1}$ es no singular y 
        \[ 
            (A^{-1})^{-1} = A.
        \]
    \item[$b)$] 
        Si $A$ y $B$ son matrices no singulares, entonces $AB$ es no singular y 
        \[ 
            (AB)^{-1} = B^{-1}A^{-1}.
        \]
    \item[$c)$] 
        Si $A$ es una matriz no singular, entonces
        \[
            {(A^\intercal)}^{-1} = {(A^{-1})}^\intercal.
        \]
    \end{enumerate}
\end{teo}

\begin{comment}
\begin{proof}\hspace{0pt} 
Demostremos el primer enunciado del teorema \ref{Sem03Clas02PropInv}. Supongamos que $A\in\Mat[\K]{n}{n}$ es una matriz no singular, entonces $A^{- 1}$ existe, además,
\[
    A^{- 1}A = AA^{- 1} = I_{n},
\]
luego, $A^{- 1}$ es no singular con inversa $(A^{- 1})^{- 1} = A$.\qedhere
\end{proof}
\end{comment}

\begin{ejem}\hspace{0pt} 
Sean $A, X\in\Mat[\K]{n}{n}$ matrices no singulares. Si $AA^{\intercal}X^{- 1} = A$ entonces $X = A^{\intercal}$.

\begin{proof}[Solución.]\hspace{0pt} Supongamos que $A$ y $X$ son matrices no singulares y que $AA^{\intercal}X^{- 1} = A$, vamos a mostrar que $X = A^{\intercal}$. Se tiene que
\begin{align*}
    AA^{\intercal}X^{- 1} = A
    &\Rightarrow A^{- 1}(AA^{\intercal}X^{- 1}) = A^{- 1}A
    && \text{Multiplicando por $A^{- 1}$ en}\\
    &\phantom{\Rightarrow A^{- 1}(AA^{\intercal}X^{- 1}) = A^{- 1}A}
    && \text{ambos lados}\\
    &\Rightarrow (A^{- 1}A)(A^{\intercal}X^{- 1}) = A^{- 1}A
    && \text{Propiedad asociativa}\\
    &\Rightarrow A^{\intercal}X^{- 1} = I_{n}
    && \text{Definición de matriz inversa}\\
    &\Rightarrow ((A^{\intercal})^{-1}A^{\intercal})X^{- 1} = (A^{\intercal})^{-1}I_{n}
    && \text{Multiplicando por $(A^{\intercal})^{-1}$ en}\\
    &\phantom{\Rightarrow ((A^{\intercal})^{-1}A^{\intercal})X^{- 1} = (A^{\intercal})^{-1}I_{n}}
    && \text{ambos lados y propiedad asociativa}\\
    &\Rightarrow X^{- 1} = (A^{\intercal})^{-1}
    && \text{Definición de matriz inversa}\\
    &\Rightarrow X = A^{\intercal}
    && \text{Teorema \ref{Sem03Clas02PropInv}}\qedhere
\end{align*}
\end{proof}
\end{ejem}

\begin{ejer}\hspace{0pt} 
Sean $A, X\in\Mat[\K]{n}{n}$ matrices no singulares tales que $(A^{\intercal})^{-1} + X^{- 1}$ es no singular. Si $2((A^{\intercal})^{-1} + X^{- 1})^{- 1} = A^{\intercal}$, entonces $X = A^{\intercal}$.
\end{ejer}

\begin{teo}
    Sean $p\in\N^*$ y $A_1, A_2, \ldots, A_p \in \Mat[\K]{n}{n}$ matrices no singulares. Se tiene que $A_1 A_2 \cdots A_p$ es no singular y 
    \[
        (A_1 A_2 \cdots A_p)^{-1} = A_p^{-1} A_{p-1}^{-1} \cdots A_1^{-1}.
    \]
\end{teo}

\begin{comment}
\begin{proof} La demostración se realiza usando Inducción, y se deja como ejercicio, el caso $p=2$ es la parte $b)$ del teorema \ref{Sem03Clas02PropInv}\qedhere
\end{proof}
\end{comment}

\begin{teo}\label{Sem03Clas01CarcInv}
    Sean $A, B \in \Mat[\K]{n}{n}$. Se tiene que si $AB=I_n$, entonces $BA=I_n$. Se dice que $B$ es la inversa de $A$ y viceversa.
\end{teo}

\begin{ejem}\hspace{0pt} 
Sean $A,B\in\Mat[\K]{n}{n}$ tales que $AB = I_{n}$, entonces $B^{2}A^{2} = I_{n}$.

\begin{proof}[Solución.]\hspace{0pt} Supongamos que $A,B\in\Mat[\K]{n}{n}$ son tales que $AB = I_{n}$, se tiene que
\begin{align*}
    B^{2}A^{2}
    &= (BB)(AA)
    && \text{Definición de $A$ a la $r$}\\
    &= B((BA)A)
    && \text{Propiedad asociativa}\\
    &= B(I_{n}A)
    && \text{Hipótesis sobre $AB$ y teorema \ref{Sem03Clas01CarcInv}}\\
    &= BA
    && \text{Multiplicar por $I_{n}$}\\
    &= I_{n}
    && \text{Hipótesis sobre $AB$ y teorema \ref{Sem03Clas01CarcInv}}\qedhere
\end{align*}
\end{proof}
\end{ejem}

\begin{ejer}\hspace{0pt} 
Sean $A,B\in\Mat[\K]{n}{n}$ tales que $AB = I_{n}$, entonces $B^{3}A^{3} = I_{n}$.
\end{ejer}

\begin{teo}
    Sea $A \in \Mat[\K]{n}{n}$, se tiene que $A$ es no singular si y solo si es equivalente por filas a $I_n$. Además,
    \[
        (A|I_n) \sim (I_n|A^{-1}).
    \]
\end{teo}
\begin{ejem} 
Dada la matriz $A = \begin{pmatrix}2&1\\1&1\end{pmatrix}$, su inversa es $A^{- 1} = \begin{pmatrix}1&-1\\-1&2\end{pmatrix}$.

\begin{proof}[Solución.]\hspace{0pt} Supongamos que $A^{- 1}$ existe, escribamos $A$ seguido de $I_{2}$ en la forma de matriz aumentada
\[
    (A|I_{2}) =
    \begin{pmatrix}
    2&1&|&1&0\\
    1&1&|&0&1
    \end{pmatrix}
\]
y después realizamos operaciones elementales por fila.
\begin{align*}
    \begin{pmatrix}
    2&1&|&1&0\\
    1&1&|&0&1
    \end{pmatrix}
    &\sim
    \begin{pmatrix}
    1&1&|&0&1\\
    2&1&|&1&0
    \end{pmatrix}
    && F_{1} \leftrightarrow F_{2}\\
    &\sim
    \begin{pmatrix}
    1&1&|&0&1\\
    0&-1&|&1&-2
    \end{pmatrix}
    &&-2F_{1} + F_{2} \to F_{2}\\
    &\sim
    \begin{pmatrix}
    1&1&|&0&1\\
    0&1&|&-1&2
    \end{pmatrix}
    &&-F_{2} \to F_{2}\\
    &\sim
    \begin{pmatrix}
    1&0&|&1&-1\\
    0&1&|&-1&2
    \end{pmatrix}
    &&-F_{2} + F_{1}\to F_{1}
\end{align*}
Dado que $A$ se redujo a $I_{2}$ se tiene que
\[
    A^{- 1} = 
    \begin{pmatrix}
    1&-1\\
    -1&2
    \end{pmatrix}.\qedhere
\]
\end{proof}
\end{ejem}

\begin{ejer}\hspace{0pt} 
Dada la matriz $A = \begin{pmatrix}1&2&0\\0&1&2\\0&0&1\end{pmatrix}$. Calcule $A^{- 1}$ si existe.
\end{ejer}


\end{document}